
\section{Thiết kế giao diện người dùng}

\subsection{Tổng quan về giao diện}

Giao diện người dùng (User Interface) là phần mà người chơi tương tác trực tiếp với chương trình. Giao diện tốt phải:
\begin{itemize}
    \item Rõ ràng và dễ hiểu.
    \item Cung cấp feedback rõ ràng cho mỗi hành động.
    \item Hướng dẫn người chơi cách nhập liệu đúng.
    \item Tránh nhầm lẫn giữa các chế độ chơi.
\end{itemize}

\subsection{Menu chính}

Menu chính là màn hình đầu tiên người dùng thấy. Nó cho phép chọn giữa ba tùy chọn:

\subsection{Hiển thị menu}

\begin{verbatim}
===== BATTLESHIP GAME =====
1. Play Game
2. Instructions
3. Quit

Choose an option (1-3): _
\end{verbatim}

\subsection{Xử lý lựa chọn}

\begin{verbatim}
IF choice == 1:
  CALL choose_game_mode()
  (Vào lựa chọn chế độ chơi)

IF choice == 2:
  CALL show_help()
  (Hiển thị hướng dẫn)

IF choice == 3:
  CALL exit_program()
  (Thoát)
\end{verbatim}

\subsection{Hướng dẫn (Help)}

Hướng dẫn hiển thị luật chơi và cách nhập liệu:

\begin{verbatim}
===== BATTLESHIP GAME - INSTRUCTIONS =====

Objective: Destroy all opponent's ships.

Rules:
1. Board size: 7x7 (coordinates 0-6)
2. Each player has:
   - 3 ships of length 2
   - 2 ships of length 3
   - 1 ship of length 4

Setup Phase:
- Place your ships on your board
- Each ship is defined by bow and stern coordinates:
  Example: "0 0 0 1" means row_bow=0, col_bow=0, row_stern=0, col_stern=1
- Ship must be horizontal (same row) or vertical (same column)

Battle Phase:
- Take turns attacking opponent's board
- Enter target coordinates: "row col"
- HIT: You hit a ship
- MISS: You hit empty water
- First to destroy all opponent's ships wins!

Press any key to continue...
\end{verbatim}

\subsection{Lựa chọn chế độ chơi}

Sau khi chọn Play Game, người chơi lựa chọn chế độ:

\subsection{Hiển thị lựa chọn chế độ}

\begin{verbatim}
===== CHOOSE GAME MODE =====
1. One Player (vs Computer)
2. Two Players

Choose a mode (1-2): _
\end{verbatim}

\subsection{Xử lý chế độ}

\begin{verbatim}
IF mode == 1:
  CALL setup_player(1)
  CALL load_auto_board()  (Tàu Computer tự động)
  gameMode = 1

IF mode == 2:
  CALL setup_player(1)
  CALL setup_player(2)
  gameMode = 2

CALL game_loop()
\end{verbatim}

\subsection{Giai đoạn Setup}

Trong giai đoạn setup, người chơi đặt tàu:

\subsection{Hướng dẫn đặt tàu}

\begin{verbatim}
===== PLAYER 1 - SETUP PHASE =====
Place your 6 ships on the 7x7 board.

Enter coordinates as: row_bow col_bow row_stern col_stern

Ship 1 (length 2): _
\end{verbatim}

\subsection{Phản hồi sau khi đặt tàu}

\begin{verbatim}
Ship placed successfully!

Ship 2 (length 2): _
\end{verbatim}

Hoặc nếu sai:

\begin{verbatim}
Invalid ship placement. Please try again.
Reason: Ship overlaps with existing ship.

Ship 1 (length 2): _
\end{verbatim}

\subsection{Hiển thị bàn cờ}

Bàn cờ được in dưới dạng lưới 7×7 với tiêu đề hàng và cột:

\subsection{Bàn cờ của chính mình}

\begin{verbatim}
YOUR BOARD - Player 1
    0 1 2 3 4 5 6
  0 [1][1][.][.][.][.][.]
  1 [.][.][.][.][.][.][.]
  2 [1][.][.][1][1][1][.]
  3 [1][.][.][.][.][.][.]
  4 [1][.][1][.][.][.][.]
  5 [.][.][1][.][.][.][.]
  6 [.][.][1][1][1][1][.]
\end{verbatim}

Legend:
\begin{itemize}
    \item [1]: Ship cell (still intact)
    \item [.]: Empty cell
    \item [X]: Destroyed ship cell
    \item [O]: Missed shot (on own board)
\end{itemize}

\subsection{Bản đồ bắn đối phương}

\begin{verbatim}
ENEMY MAP - Player 2
    0 1 2 3 4 5 6
  0 [.][X][.][.][.][.][.]
  1 [O][.][.][.][.][.][.]
  2 [.][.][.][.][.][.][.]
  3 [.][.][.][.][.][.][.]
  4 [.][.][X][.][.][.][.]
  5 [.][.][.][.][.][.][.]
  6 [O][.][.][.][.][.][.]
\end{verbatim}

Legend:
\begin{itemize}
    \item [.]: Unknown cell (not attacked)
    \item [O]: Miss
    \item [X]: Hit
\end{itemize}

\subsection{Dual-Board Display}

Trong giai đoạn chiến đấu, hai bàn cờ được in song song để người chơi thấy cả hai cùng lúc:

\begin{verbatim}
YOUR BOARD - Player 1          ENEMY MAP - Player 2
    0 1 2 3 4 5 6                  0 1 2 3 4 5 6
  0 [1][1][.][.][.][.][.]       0 [.][X][.][.][.][.][.]
  1 [.][.][.][.][.][.][.]       1 [O][.][.][.][.][.][.]
  2 [1][.][.][1][1][1][.]       2 [.][.][.][.][.][.][.]
  3 [1][.][.][.][.][.][.]       3 [.][.][.][.][.][.][.]
  4 [1][.][1][.][.][.][.]       4 [.][.][X][.][.][.][.]
  5 [.][.][1][.][.][.][.]       5 [.][.][.][.][.][.][.]
  6 [.][.][1][1][1][1][.]       6 [O][.][.][.][.][.][.]
\end{verbatim}

\subsection{Giai đoạn chiến đấu}

Trong giai đoạn chiến đấu:

\subsection{Yêu cầu nhập tọa độ}

\begin{verbatim}
===== PLAYER 1 TURN =====
Enter target row and column (e.g., "2 3"): _
\end{verbatim}

\subsection{Thông báo kết quả}

\begin{verbatim}
HIT!
(Cập nhật bàn cờ và bản đồ bắn)
\end{verbatim}

Hoặc:

\begin{verbatim}
MISS!
(Cập nhật bản đồ bắn)
\end{verbatim}

Hoặc nếu input sai:

\begin{verbatim}
You already attacked this cell! Try another.

Enter target row and column (e.g., "2 3"): _
\end{verbatim}

\subsection{Tạm dừng để đổi lượt (Pause for Swap)}

Trong chế độ Two Players, sau khi mỗi lượt, chương trình tạm dừng để người chơi khác có cơ hội giảm bớt việc nhìn thấy bàn cờ của đối thủ:

\begin{verbatim}
===== PLEASE SWAP PLAYERS =====
Player 2, it's your turn now.
(Press any key to continue...)
\end{verbatim}

Thủ tục pause\_for\_swap:

\begin{verbatim}
PROCEDURE pause_for_swap()
  PRINT message
  READ one character (to wait for user confirmation)
  
  (Optionally, clear screen to hide previous board)

END PROCEDURE
\end{verbatim}

\subsection{Thông báo kết thúc game}

Khi một người chơi chiến thắng:

\begin{verbatim}
===== GAME OVER =====
PLAYER 1 WINS!

Congratulations! You destroyed all opponent's ships.

Return to menu? (Y/N): _
\end{verbatim}

Hoặc:

\begin{verbatim}
===== GAME OVER =====
COMPUTER WINS!

Game Over. Better luck next time!

Return to menu? (Y/N): _
\end{verbatim}

\subsection{Thông báo lỗi}

Chương trình cung cấp thông báo lỗi rõ ràng:

\begin{enumerate}
    \item \textbf{Invalid input format}: Khi không thể tách số từ chuỗi input.
    \item \textbf{Coordinates out of range}: Khi tọa độ ngoài [0, 6].
    \item \textbf{Ship overlaps with existing ship}: Khi đặt tàu chồng lấn.
    \item \textbf{Ship must be horizontal or vertical}: Khi tàu chéo.
    \item \textbf{Ship length does not match}: Khi độ dài sai.
    \item \textbf{You already attacked this cell}: Khi bắn lại cùng ô.
\end{enumerate}

\subsection{Nguyên tắc thiết kế}

Giao diện được thiết kế theo các nguyên tắc:

\begin{itemize}
    \item \textbf{Clarity}: Rõ ràng về những gì người chơi cần làm tiếp theo.
    \item \textbf{Consistency}: Cùng format, cùng ký hiệu ở mọi nơi.
    \item \textbf{Feedback}: Phản hồi rõ ràng cho mỗi hành động.
    \item \textbf{Error Handling}: Thông báo lỗi cụ thể, không chung chung.
    \item \textbf{Accessibility}: Dễ sử dụng, không cần huấn luyện phức tạp.
\end{itemize}
